\section{Falsifiability}

The claims of this paper are falsifiable within the OC/ETS framework. We state
explicit structural conditions under which the phase-transition interpretation
of enterprise identity loss would be invalid.

\paragraph{Falsification via Cycle Persistence.}
If, after an alleged identity loss event, there exists a non-empty set of cycles
$C_I \subset C_{\mathrm{EA}}$ whose trajectories remain entirely within an
admissible subspace $\Lambda_I \subset \Omega_{\mathrm{EA}}$ that is bounded away
from $\partial\Omega_{\mathrm{EA}}$, then identity loss has not occurred in the
formal sense defined here. Observation of such persistent invariant cycles
falsifies the identity-loss claim.

\paragraph{Falsification via Admissible Reconstitution without Extension.}
If identity-supporting cycles can be regenerated under the closed dynamics of
$\mathrm{ETS}.K_{\mathrm{EA}}$ without introducing new admissible structure (no
new axes, potentials, thresholds, or state variables), then the irreversibility
and no-go results presented in this work are false.

\paragraph{Falsification via Absence of Discontinuity.}
If the transition attributed to identity loss can be continuously deformed into
ordinary operational fluctuation within $\Omega_{\mathrm{EA}}$ without boundary
attachment, cycle extinction, or threshold crossing, then the phenomenon does
not constitute a phase transition in the OC sense.

\paragraph{Falsification via Collapse Equivalence.}
If every empirically admissible instance of identity loss necessarily coincides
with collapse of $\Omega_{\mathrm{EA}}$ or with $k_{\mathrm{EA}} \rightarrow 0$,
then identity loss is not structurally distinct from organizational failure.
Such a finding would falsify the central distinction advanced in this paper.

These falsifiability conditions ensure that identity loss, as formalized here,
is neither vacuous nor purely interpretive, but a testable structural claim
within the OC/ETS framework.
